%% kapa: exam=1 round=37 year=2026 period=2 form=A subject=accounting subject_ko=회계학 seq=39 number=79
%% source: Q-Net 감정평가사 기출문제 (cst00309) artlSeq=5258924 file=raw/1008/5258924/2026년 제37회 감정평가사 1차 2교시.pdf
%% conversion: kordoc 4.24.51 -> tools/gen1.py (auto); review=proofread-sonnet; flags=none
\begin{question}{79}
(주)감평은 정상원가계산제도를 채택하고 있으며, 제조간접원가는 직접노무원가의 50 \%를 예정배부한다. (주)감평은 선입선출법에 의해 재고자산을 평가하고 있으며, 당기 기초 및 기말 재고자산과 원가 관련 자료는 다음과 같다.
\par
{\small\centering
\begin{tabular}{@{}lrr@{}}\hline
 & \multicolumn{1}{c}{기초재고} & \multicolumn{1}{c}{기말재고} \\\cline{2-3}
직접재료 & ￦12,000 & ￦15,000 \\
재 공 품 & 0 & ? \\
제 \quad 품 & 48,000 & 17,000 \\
\multicolumn{3}{@{}l}{\hspace*{1em}○ 직접재료 매입액은 ￦120,000이다.} \\
\multicolumn{3}{@{}l}{\hspace*{1em}○ 직접노무원가 발생액은 기본원가의 40 \%이다.} \\\hline
\end{tabular}\par}
당기에 발생한 직접재료원가 중 ￦99,000이 당기제품제조원가에 포함되어 있으며, 기말재공품에 예정배부된 제조간접원가가 ￦6,000일 경우, 기말재공품의 원가는? (단, 공손 및 감손은 발생하지 않는다.) \par
\begin{choices}
  \item ￦31,400
  \item ￦34,600
  \item ￦36,000
  \item ￦37,600
  \item ￦40,200
\end{choices}
\end{question}
\answer{3}
